% =========================================================
% INTRODUÇÃO
% FINALIDADE: situar o leitor no tema e no problema da pesquisa.
% ESCREVA: contexto, problema, justificativa, objetivos e organização do
%   trabalho, no lugar dos parágrafos de orientação.
% CUIDADO: não antecipe resultados. As subseções são sugestão e podem ser
%   ajustadas com a orientação.
% EXEMPLOS NESTE ARQUIVO: seção, subseção, lista de objetivos e remissão.
% =========================================================

% \chapter = seção primária (1); \section = secundária (1.1);
% \subsection = terciária (1.1.1); \subsubsection = quaternária (1.1.1.1).
\chapter{INTRODUÇÃO}
\label{cap:introducao}

A introdução é a parte inicial do texto, na qual se apresentam o tema, a delimitação do problema de pesquisa, a justificativa, os objetivos e a organização da dissertação. Deve situar o leitor quanto ao assunto tratado e ao ponto de vista adotado, sem antecipar resultados ou conclusões.

\section{Contextualização e problema de pesquisa}
\label{sec:problema}

% Apresente o tema, o contexto em que se insere e o recorte
% adotado (espacial, temporal, população ou corpus). Conclua
% com a formulação do problema, preferencialmente em forma
% de pergunta.
Nesta seção, apresenta-se o tema da pesquisa e o contexto em que ele se insere, com apoio na literatura da área. A exposição parte de um quadro mais amplo e avança até o recorte específico do estudo, de modo a evidenciar a lacuna ou a situação que motivou a investigação. Ao final, formula-se o problema de pesquisa, de preferência na forma de uma pergunta clara e delimitada.

\section{Justificativa}
\label{sec:justificativa}

% Indique a relevância científica, social e educacional do
% estudo e sua relação com a linha de pesquisa do programa.
A justificativa expõe as razões que tornam a pesquisa relevante. Convém indicar a contribuição esperada para o conhecimento da área, para a prática educativa e para a sociedade, bem como a relação do estudo com a linha de pesquisa do Programa e, quando pertinente, com a trajetória acadêmica e profissional do(a) autor(a).

\section{Objetivos}
\label{sec:objetivos}

\subsection{Objetivo geral}

% Em geral, um objetivo amplo, iniciado por verbo no infinitivo e
% coerente com o problema de pesquisa.
Analisar [objeto de estudo] em [contexto], com vistas a [finalidade da pesquisa].

\subsection{Objetivos específicos}

% A quantidade de objetivos específicos depende das necessidades da
% pesquisa e deve ser definida com o(a) orientador(a). Inicie cada um
% por verbo no infinitivo. Cada objetivo deve corresponder a uma etapa
% verificável da pesquisa e ser retomado nos resultados e nas
% considerações finais. Para acrescentar ou remover, copie ou apague
% uma linha \item.
\begin{enumerate}[label=\alph*), leftmargin=1.5cm+\labelwidth, itemsep=0pt]
  \item identificar [primeiro aspecto do objeto];
  \item descrever [segundo aspecto do objeto];
  \item analisar [terceiro aspecto do objeto];
  \item propor [produto ou encaminhamento, quando houver].
\end{enumerate}

\section{Estrutura do trabalho}
\label{sec:estrutura}

% Descreva brevemente o conteúdo de cada seção. As remissões
% abaixo são atualizadas automaticamente pelo LaTeX.
% Remissão: \label{nome} dá um nome à seção (veja cada \chapter);
% \ref{nome} escreve o número dela. O ~ impede a quebra de linha.
Além desta introdução, a dissertação está organizada em mais quatro seções. A Seção~\ref{cap:fundamentacao} apresenta a fundamentação teórica que orienta a pesquisa. A Seção~\ref{cap:metodologia} descreve o percurso metodológico. A Seção~\ref{cap:resultados} apresenta e discute os resultados obtidos. A Seção~\ref{cap:consideracoes} reúne as considerações finais, as limitações do estudo e as indicações para pesquisas futuras.
